\section{LL(k) versus LR(k) dan Pemilihan Strategi}
\label{sec:ch12-ll-lr}

\subsection{Perbandingan}

\begin{table}[H]
\centering
\caption{Perbandingan keluarga LL(k) dan LR(k)}
\label{tab:ch12-ll-lr}
\begin{tabularx}{\textwidth}{@{}lXX@{}}
\toprule
\textbf{Aspek} & \textbf{LL(k)} & \textbf{LR(k)} \\
\midrule
Arah & Kiri ke kanan, derivasi kiri & Kiri ke kanan, derivasi kanan \\
Lookahead & \(k\) token & \(k\) token \\
Left recursion & Tidak didukung & Didukung \\
Grammar & Harus difaktorkan & Lebih alami \\
Tabel & Relatif kecil & Lebih besar \\
Kesalahan & Mudah dilokalisasi & Perlu strategi tambahan \\
Implementasi & Recursive descent & Stack dan tabel \\
\bottomrule
\end{tabularx}
\end{table}

\subsection{Pemilihan}

LL(k) cocok untuk:
\begin{itemize}[leftmargin=*]
  \item parser tulisan tangan;
  \item grammar yang dapat difaktorkan;
  \item prioritas pada pesan kesalahan.
\end{itemize}

LR(k) cocok untuk:
\begin{itemize}[leftmargin=*]
  \item grammar besar dan kompleks;
  \item parser generator;
  \item left recursion alami;
  \item cakupan bahasa yang lebih luas.
\end{itemize}

\subsection{Di Balik Layar Bison}

Bison membangun tabel LALR(1) dari aturan grammar. Pengguna menulis
aturan, deklarasi token, dan semantic action. Bison menghasilkan
fungsi \code{yyparse} yang memakai tabel tersebut. Jika terjadi
konflik, Bison melaporkan jumlah dan lokasinya.

Pada Bab XIII, kita akan menulis grammar Bison untuk Subset-C,
menambahkan semantic action, dan membangun AST. Pemahaman item dan
konflik pada bab ini membantu membaca laporan Bison dan memperbaiki
grammar.

\begin{infobox}
Tidak ada strategi yang selalu terbaik. Pilihan bergantung pada
ukuran bahasa, kebutuhan diagnostik, ketersediaan alat, dan
keterampilan tim. Yang penting adalah memahami trade-off dan
mendokumentasikan keputusan.
\end{infobox}
